\section[第7章\q 打肿脸充胖子]{第7章\\ 打肿脸充胖子}
\label{sec:ch7}

“拜伦，”一名娇小优雅的女子走到拜伦旁边，拍着他的胳膊，“你没为难人家小妹妹吧？”

“柔伊，”拜伦说，“这是我太太，叫她乔琪雅就好。这位是亨利的朋友，柔伊。”拜伦转向柔伊，小声嘀咕，虽然这音量在五张桌子之外都听得到，“我就是不懂。她总是站在这里，看着同样的照片。好像不盯着看，照片就会发生什么变化似的。”

拜伦的妻子不理他，迳自走去亨利之前坐着的椅子上。“拜伦来就是找人聊天的，”她说，“我来是看画的。很高兴认识你，柔伊。”她伸手来，优雅地握住了柔伊的手。

柔伊笑着说，“我来又聊天又看画，还要享用咖啡。”

“是的。”乔琪雅同意。拜伦开始吃栉瓜派时，她吹着她的浓缩咖啡。“这个海伦娜咖啡，还真不同凡响。”

“如果你们不介意，”柔伊说。“我想请教，你们怎么认识亨利的？”

\subsection{累积财富不容易，要花掉却很容易}

“哇，那是个很有意思的故事。”拜伦抢在妻子前头说。“15年应该有啰？”

“18年啦！”乔琪雅说。

拜伦耸肩。“对，18年。还真是有点久了。那时景气之好，石油产业当红炸子鸡。是我们的黄金时代。我是个得意的生意人。奥克拉荷马的石油之王。富可敌国。但是其实我就是一个白痴。真正的生意头脑在这儿，”他歪头示意乔琪雅的方向。“那时的我太过冥顽不灵，哪里会知道。”说完，又对着妻子说：“对吗？”

“不予置评。”乔琪雅说。

“我们德州有一句话，叫做大帽子底下，没有牛（Big Hat, not Cattle），意思是打肿脸充胖子，”拜伦继续说道，“都是表面的，没有实业做支撑。那就是我。我操作太过。”他又朝乔琪雅点点头，乔琪雅只管喝她的浓缩咖啡。“因为我把她也捆住了，所以，可以说我们操作太过。景气突然改变了。石油产业不再风光一时，不再遍地黄金。”

他又看了看柜台后方忙着准备客人咖啡的亨利。“他对你说过，怎样累积财富吧？”

“一次存一点。”柔伊回答。

乔琪雅说：“事实是，亏损起来，也是同样的道理。”

“是呀！”拜伦说，回看着柔伊。“人生第一个100万，你知道的，”他摇头说，“\textbf{累积起来不容易，要花掉却容易极了}。还没意识到，早就花光光了。还不弄明白怎么花的，”他耸着肩膀说，“就已经没了。

“所以有一天，我穿着高领紧身衫，坐在医师诊所里，而医生忙着调整机器，就是听诊器之类的，听着我的心跳，然后清了清喉咙说：‘拜伦，我要和你说几句话，你用心听好啦？’然后，他就说，如果我继续喝酒、抽烟，一天吃很多肉，过不久就会死。

“我回答他说，‘直话直说吧！告诉我，我到底怎么了？’

“‘那我就直说了。拜伦，自己做一个选择吧！要继续保有各种恶习，还是要命？’

“于是我就坐在诊疗室的椅子上，足足看了他1分钟。他忍不住开口问：‘怎么样？你也开口说说嘛！’

“再给我一点时间，我总得想一想。”

他仰头大笑，在前面柜台的客人回头看过来，不晓得发生了什么事。而坐在附近的客人忍不住笑了，却不往这边看。柔伊想，他们大概是常客，所以听过拜伦的故事。

“我也真的做了一番思考，而且是用了几个月的时间。也不是什么都不做的坐在那里想，就是确实把这件事放在心上。大部分时候的想法是，你说什么？不吃肉？不抽烟？滴酒不沾？别和我说这些了。并不是因为我酗酒而有问题，而是什么都有问题。超重36公斤，又过度自大。好像在比赛着，看是我的信用先破产，还是婚姻先破产。”

“还有你的心血管。”乔琪雅说。

“对了，忘记说那个了。”

“还有你可能忘了，还有你的心脏绕道手术。”

拜伦呵呵笑着说。“是的，那个也忘了。”

柔伊喜欢看他们两个斗嘴，好像一对喜剧演员。这两位年纪都足以做她的长辈了。上回和父母碰面是什么时候？

“拜伦因公出差来纽约时，”乔琪雅说，“他心脏病发作。我从土尔沙赶到时，他已经要进手术房了。”

“亲爱的，”10秒钟没说话，拜伦已然开始焦虑了。“我正要说到那件事。柔伊，你想再一杯咖啡吗？”

“谢谢。不用。”柔伊说。

拜伦朝柜台大步走去，柔伊转身对着乔琪雅。“所以他已经要进行手术了。”她提示乔琪雅讲下去。

乔琪雅一下说不出话来。“我们——”柔伊突然发现女人眼睛已经满是眼泪，“我们差点就要失去他了，他却顽固得好像一只乌龟。”乔琪雅笑着，用一小块餐巾纸拭擦着眼角。“谁的话都不听，不管是医师还是老婆说的，连女儿说的都没用。”乔琪雅长长呼了一口气，喝了一口特浓咖啡。“但他听亨利的。”

“听亨利的？”柔伊想像着奥克拉荷马的石油大亨和布鲁克林区咖啡店早班组长狭路相逢的情景。

“是的。”乔琪雅说，“我们不能回去，也不能在短时间内旅行。即便在拜伦出院后，我们还是只能在此逗留一阵子。就像电影台词说的‘别想试图离开这里’。”她哑然失笑，“我们去逛布鲁克林区各式的艺术展览。有一天，来这里喝咖啡时，被墙上的照片吸引。就是这样遇见亨利的，他和拜伦一见如故，亨利开始对他说教。”

“如何理财？”

乔琪雅笑着说：“关于怎么吃，怎么生活，怎样活得好。我记得第一回见面时，亨利说过‘创建健康的过程和累积财富的过程是一样的，拜伦’。”

两个女人异口同声的说道：“一次存一点。”

乔琪雅脸上的线条变成柔顺了，“拜伦还真的听进去了。”她摇头，“感谢老天给了我们一个奇迹。但是，很快的又出现了财务问题。特别是收到医院帐单后，换我要得心脏病了。”她看着柔伊说：“我好像听见亨利说‘其实你已经很富有了，只是你不知道？’”

柔伊点头，心里想说，哇，这个女人什么都知道。

“对我们来说，刚好相反。我们已经破产了，但是我们却不知道。当然，那时知道了。”

她停了一下，看着柔伊。

“我不知道我们透支的情形有多严重。我从来不问钱的事，都是拜伦在管，直到他躺到医院里，我才打开了邮件，发现真相。我们身上几乎没有余钱。所有房屋都贷了许多款，信用卡债一大堆，我们只能支付每月最低应缴金额。”

柔伊心里抖了一下。每个月付最低应缴金额，不就是柔伊本人处理信用卡的方法吗？

“我了解问题很严重。”乔琪雅继续说，“我不懂的是，他怎么走到一步的。有一天亨利解释给我听，他说‘如果你欠了2万美元的信用卡债，每月只付最低应缴金额，要多花18年才会还完，期间支付的总金额超过4.6万美元。’

“我听完脸都白了。那比我们欠款的2倍还多。”

“喔！”柔伊回答。

“那还只是一张信用卡。”乔琪雅说，“要知道我们财务上的洞有多大，可能要把这个数字放大许多倍，复利的神奇力量。我猜，你已经听亨利说过吧？”

柔伊点头。

“这是一把双面刃。可以帮你，也可以毁了你。\textbf{欠债也是利滚利的}，而且一旦开始，就会迅速成长，且一发不可收拾。”乔琪雅摇头，“我还以为我们过得不错。结果完全不是。我们把自己的脸打得很肿。”

柔伊突然想起了周一早上看到的那只广告。

{\kaishu 如果你不知道要去哪里，只怕到了哪都不会快乐。}

\subsection{问题并非缺钱，而是缺少沟通与方法}

“后来，你们怎么办？”

“是这样的，”乔琪雅说，“亨利帮助我们慢慢走过来。我们先卖掉在土尔沙的一栋商业大楼，两处度假屋。卖掉之后，扣掉贷款，刚好够我们在曼哈顿区买一处公寓。我们之后没有再回奥克拉荷马，然后再一点点重建生活。”

“拜伦开始戒烟后，我们发现抽烟很花钱。不单是健康问题，还有每个月多出的一项开支。”她笑了，“你听过亨利说他的拿铁因子吗？”

柔伊点头，并在脑中划下一个重点，这个稍后一定要再问一下亨利。因为到现在，她还不知道拿铁因子的具体意思。

“总之，”乔琪雅继续说，“拜伦叫它是‘香烟因子’。他有时候也叫它为雪茄因子。”她用手在脸前头晃了晃，好像要把浓厚的烟味驱走。

“把这些都戒掉之后，不到1个月的时间，他好像换了个人似的。那些咳嗽、病痛都不见了。但我们的财务上还有个大洞。”

“我们停用信用卡，花了几年还清卡债，但我们还是还完了。之后，买了二手车。我老妈常说：‘船到桥头自然直。’关于钱不够用的问题，她有句箴言：‘买二手货，付现金。’”

她把最后一口浓缩咖啡喝完，并把杯子倒扣在桌上。

“亨利不单拯救了我们的财务危机、拜伦的健康，还有我们的婚姻。因为我们谈到钱的事情时，不是争吵，而是沟通，一起解决问题。”

“金钱曾让我们的关系有如骨肉分离，现在金钱又把我们黏在一起。”

“对了，你结婚了吗？”乔琪雅问， 柔伊摇头，“记住，金钱是导致离婚的第一大元凶，但不是因为金钱本身，也不是因为缺少金钱，\textbf{而是因为缺少如何使用金钱的沟通，和处理与钱相关事物的方法}。”

“我永远都不会忘记，我们有天真正的坐下来，第一次敞开来讨论钱的事情，讨论我们的生活和未来。讨论我们真正需要什么，如何来达成。那一回，我们都打破了沉默。坐在厨房餐桌旁边，开诚布公。”

乔琪雅笑了。

“就我所知，我还是世上最富有的人。”

他们沉默了一会儿。

“就是这样，我们后来常常来亨利这里。不久之后，亨利就在拜伦工作的能源公司投入了一笔不小的金额，两个男人就整天讨论什么新技术，什么新显学……怎么了，你还好吧？”

柔伊听完，脑中轰的一声。一笔不小的金额？亨利投资？她的脑子在飞转。“等下。你说，之后你就一直来亨利的咖啡厅？意思是，他是这里的经理？”

乔琪雅伸出手来，拍了拍柔伊的手。“亲爱的小妹妹，亨利是这儿的老板，这家店是他从无到有创建起来的。”

“乔琪雅，”是拜伦，在她们后方，轻敲着桌面。

“好了，”乔琪雅站起身来说。“我们要去机场和女儿碰面了。下周见，亨利。”他们一边往门口走，一边说道：“很高兴认识你，柔伊。”

柔伊看着拿铁发呆，头也来不及抬起来。
